% ============================================================================
% Back Matter — Glossary
% ============================================================================

\clearpage
\phantomsection
\addcontentsline{toc}{chapter}{Glossary}

% ---- Manual Glossary (styled as a chapter) ----
\chapter*{Glossary}
\markboth{Glossary}{Glossary}

\begin{description}[leftmargin=2cm, style=nextline, font=\bfseries\sffamily\color{NavyBlue}]

\item[ALU (Arithmetic Logic Unit)]
The part of the CPU that performs all arithmetic (maths) and logical (comparison) operations.

\item[Application Software]
Programs designed for end-users to perform specific tasks, such as MS Word, Chrome, or Calculator.

\item[AutoFill]
An Excel feature that automatically fills cells with a series or pattern based on initial values.

\item[Backup]
A copy of data stored separately to protect against data loss.

\item[Boot Process]
The sequence of events that occurs when a computer is turned on, from POST to the desktop appearing.

\item[Browser]
An application used to view websites on the Internet (e.g., Chrome, Edge, Firefox).

\item[Cell Address]
The unique identifier of a cell in Excel, made of its column letter and row number (e.g., B3).

\item[Cloud Storage]
Online storage that allows files to be saved and accessed over the Internet (e.g., Google Drive, OneDrive).

\item[CPU (Central Processing Unit)]
The ``brain'' of the computer that processes instructions and manages data.

\item[Cyberbullying]
Using technology to harass, threaten, or humiliate someone online.

\item[Desktop]
The main screen of the Windows operating system where icons, taskbar, and windows appear.

\item[Device Driver]
A small program that allows the operating system to communicate with a specific hardware device.

\item[Digital Citizenship]
Responsible, ethical, and safe use of technology and the Internet.

\item[Email]
Electronic mail; a method of sending digital messages over the Internet.

\item[Ethernet]
A wired network connection using a cable plugged into a router or network port.

\item[Extension (File)]
The suffix after the dot in a filename that indicates the file type (e.g., .docx, .pdf, .jpg).

\item[File]
A collection of data stored with a name and extension on a storage device.

\item[Firmware]
Permanent software built into hardware devices (e.g., BIOS on a motherboard).

\item[Folder]
A container used to organise files on a storage device; also called a directory.

\item[Formula]
An expression in Excel that performs calculations using cell references, operators, and functions.

\item[Function (Excel)]
A pre-built formula in Excel for common calculations (e.g., SUM, AVERAGE, MAX, MIN).

\item[GUI (Graphical User Interface)]
A visual interface using windows, icons, menus, and pointers to interact with the computer.

\item[Hardware]
The physical components of a computer that you can touch (e.g., monitor, keyboard, CPU).

\item[Internet]
A global network of interconnected computers that communicate using standard protocols.

\item[IoT (Internet of Things)]
The network of everyday physical devices connected to the Internet, capable of collecting and sharing data.

\item[ISP (Internet Service Provider)]
A company that provides Internet access to homes, schools, and businesses.

\item[Keyboard Shortcut]
A combination of keys pressed together to perform a quick action (e.g., Ctrl+C to copy).

\item[Malware]
Malicious software designed to damage, disrupt, or gain unauthorised access to a computer.

\item[Operating System (OS)]
System software that manages hardware, provides a user interface, and runs application software.

\item[Phishing]
A type of online scam using fake emails or websites to trick users into sharing personal information.

\item[POST (Power-On Self-Test)]
A diagnostic test run by the computer at startup to check that essential hardware is working.

\item[RAM (Random Access Memory)]
Volatile memory used to store data temporarily while programs are running.

\item[Ribbon]
The toolbar area in MS Office applications containing tabs, groups, and commands.

\item[ROM (Read-Only Memory)]
Non-volatile memory containing permanent instructions (e.g., BIOS firmware).

\item[Search Engine]
A tool that helps find information on the Internet (e.g., Google, Bing, Yahoo).

\item[Slide (PowerPoint)]
A single page in a PowerPoint presentation that displays content.

\item[Software]
A set of instructions (programs) that tells the computer hardware what to do.

\item[Spreadsheet]
A document organised in rows and columns of cells, used for data, calculations, and charts.

\item[System Software]
Software that manages hardware and provides a platform for other software (e.g., Windows, Linux).

\item[Taskbar]
The bar at the bottom of the Windows desktop showing open programs, pinned apps, and system tray.

\item[Transition (PowerPoint)]
A visual effect that plays when moving from one slide to the next.

\item[URL (Uniform Resource Locator)]
The web address used to locate a specific page on the Internet (e.g., https://www.google.com).

\item[USB (Universal Serial Bus)]
A standard connector for attaching devices (flash drives, keyboards, etc.) to a computer.

\item[Virus]
A type of malware that spreads by infecting other files and can damage the computer system.

\item[Wi-Fi]
A wireless networking technology that allows devices to connect to the Internet without cables.

\item[World Wide Web (WWW)]
A collection of websites and web pages accessible via the Internet using a browser.

\end{description}
